\documentclass{ctexbook}
\usepackage{ctex}
\usepackage{amsmath}
\usepackage{amssymb}
\usepackage[colorlinks,linkcolor=blue]{hyperref}
\usepackage{booktabs}
\usepackage[paper=a4paper, left=1.5cm, right=1.5cm, top=3cm, bottom=3cm]{geometry}
\usepackage{enumitem}
\usepackage{tikz}
\usepackage{pgfplots}
\usepackage{amsthm}

\newcommand{\ITEMS}{\begin{enumerate}[leftmargin=3.5em]}
\newcommand{\ITEME}{\end{enumerate}}
\newcommand{\no}{\noindent}
\newcommand{\R}{\mathbb{R}}
\newcommand{\N}{\mathbb{N}}

\setcounter{tocdepth}{2}



\newtheorem{theorem}{定理}[chapter]
\newtheorem*{lemma}{引理}
\newtheorem{proposition}[theorem]{命题}
\newtheorem{corollary}[theorem]{推论}
\newtheorem*{remark}{\indent 注}

\newtheorem{example}{Example}[section]




% update list
% 
%  完成章节：1
%  
%  实数类部分
%  微分学
%
%
%
%
%
%
%
%
%




\title{数学分析}
\author{thchXR}

\begin{document}

\maketitle

\setcounter{page}{1}
\pagenumbering{Roman}

\tableofcontents

\setcounter{page}{1}
\pagenumbering{arabic}
% ================================================================================== 
% 						逻 辑 符 号 与 离 散 数 学 基 础
% ================================================================================== 
\chapter{逻辑符号与离散数学基础}

离散数学是现代数学的重要分支，也是计算机科学的核心基石，主要研究不连续的、离散的数学结构和对象（如整数、命题和图）。这里我只介绍最基本的逻辑符号，其他内容将在后续离散数学的章节详细介绍。

\section{命题逻辑}

\subsection{命题与联接词}

命题是逻辑的基本成分，命题是一个或真或假的陈述语句，但不可能即真又假。我们习惯用小写字母 $p,q,r,s,\cdots$ 表示命题。如果一个命题为真，用 T 表示，反之用 F 表示。

特别的，对于无法拆分成更简单语句的命题，我们称之为简单命题或原子命题，而称由联接词和简单命题构成的命题为复合命题。

\subsubsection{命题常项与命题变项}

对于简单命题来说，其真值是确定的，因而又称命题常项或命题变元，而对于一些含有变量的陈述句，虽然不是命题，但只要给定变量的值就可以得到其真值，所以我们称之为命题变项或命题变元。

\subsubsection{否命题}

设有一命题 $p$ ，则 $p$ 的否定记为 $\neg p$（也可记为 $\overline p$）， $p$ 的真值与 $\neg p$ 的真值相反。

\subsubsection{合取}
设 $p,q$ 为命题， $p,q$ 的合取式为复合命题，记为 $p\land q$ ，$\land$ 是合取联接词 ，即 $p,q$ 均为 T 时 $p\land q$ 取 T ，否则 $p\land q$ 取 F。

\subsubsection{析取}
设 $p,q$ 为命题， $p,q$ 的析取式为复合命题，记为 $p\lor q$ ，$\lor$ 是析取联接词 ，即 $p,q$ 均为 F 时 $p\lor q$ 取 F ，否则 $p\lor q$ 取 T。

\subsubsection{异或}
设 $p,q$ 为命题， $p,q$ 的异或为复合命题，记为 $p\oplus q$ ，$\oplus$ 是异或联接词，即 $p,q$ 的真值不同时 $p\oplus q$ 取 T ，否则 $p\oplus q$ 取 F。

\subsubsection{真值表}
对于以上命题，我们可以用以下真值表表示
\smallskip

% 真值表 1-1
\begin{minipage}[c]{0.25\textwidth}
	\centering
	\begin{tabular}{c|c}
	$p$ & $\neg p$ \\
	\hline
	T & F \\
	F & T \\
	\end{tabular}
\end{minipage}
\begin{minipage}[c]{0.25\textwidth}
	\centering
	\begin{tabular}{cc|c}
	$p$ & $q$ & $p\land q$ \\
	\hline
	T & T & T \\
	T & F & F \\
	F & T & F \\
	F & F & F \\
	\end{tabular}
\end{minipage}
\begin{minipage}[c]{0.25\textwidth}
	\centering
	\begin{tabular}{cc|c}
	$p$ & $q$ & $p\lor q$ \\
	\hline
	T & T & T \\
	T & F & T \\
	F & T & T \\
	F & F & F \\
	\end{tabular}
\end{minipage}
\begin{minipage}[c]{0.25\textwidth}
	\centering
	\begin{tabular}{cc|c}
	$p$ & $q$ & $p\oplus q$ \\
	\hline
	T & T & F \\
	T & F & T \\
	F & T & T \\
	F & F & F \\
	\end{tabular}
\end{minipage}


\subsection{条件语句}

\subsubsection{蕴涵}
设 $p,q$ 为命题，命题 $p \rightarrow q$ 为命题 “若 $p$ ，则 $q$ ”， $\rightarrow$ 是蕴涵联接词。当 $p$ 为真而 $q$ 为假时，命题 $p \rightarrow q$ 为假，否则 $p \rightarrow q$ 为真，对于命题 $p \rightarrow q$ ，其中 $p$ 为前提， $q$ 为结论。

\begin{remark}
为防止后续极限符号 $\to$ 与蕴涵符号产生歧义，在后续分析学章节中将使用 $\Rightarrow$ 符号，仅在讨论逻辑学时使用 $\rightarrow$。	
\end{remark}


\subsubsection{逆，倒置与反}
设 $p,q$ 为命题，对于条件语句 $p \rightarrow q$ ，称 $q \rightarrow p$ 为 $p \rightarrow q$ 的逆蕴涵，称 $\neg q \rightarrow \neg p$ 为 $p \rightarrow q$ 的倒置蕴涵，称 $\neg p \rightarrow \neg q$ 为 $p \rightarrow q$ 的反蕴涵。
\smallskip

% 真值表 1-2
\begin{minipage}[c]{0.25\textwidth}
	\centering
	\begin{tabular}{cc|c}
	$p$ & $q$ & $p\rightarrow q$ \\
	\hline
	T & T & T \\
	T & F & F \\
	F & T & T \\
	F & F & T \\
	\end{tabular}
\end{minipage}
\begin{minipage}[c]{0.25\textwidth}
	\centering
	\begin{tabular}{cc|c}
	$p$ & $q$ & $q \rightarrow p$ \\
	\hline
	T & T & T \\
	T & F & T \\
	F & T & F \\
	F & F & T \\
	\end{tabular}
\end{minipage}
\begin{minipage}[c]{0.25\textwidth}
	\centering
	\begin{tabular}{cc|c}
	$p$ & $q$ & $\neg q \rightarrow \neg p$ \\
	\hline
	T & T & T \\
	T & F & F \\
	F & T & T \\
	F & F & T \\
	\end{tabular}
\end{minipage}
\begin{minipage}[c]{0.25\textwidth}
	\centering
	\begin{tabular}{cc|c}
	$p$ & $q$ & $\neg p \rightarrow \neg q$ \\
	\hline
	T & T & T \\
	T & F & T \\
	F & T & F \\
	F & F & T \\
	\end{tabular}
\end{minipage}

\subsubsection{等价}

设 $p,q$ 为命题，命题 $p \leftrightarrow q$ 为命题 “$p$ 当且仅当 $q$ ”，我们也将当且仅当写作 iff （if and only if）， $\leftrightarrow$ 是等价联接词。当 $p$ 与 $q$ 的真值相等时，命题 $p \leftrightarrow q$ 为真 ，否则命题 $p \leftrightarrow q$ 为假。

\begin{remark}
在后续分析学章节中将使用 $\Leftrightarrow$ 符号表示等价，仅在讨论逻辑学时使用 $\leftrightarrow$。	
\end{remark}

\section{命题公式与等值演算}

\subsection{命题公式与分类}

\subsubsection{命题公式}

为了更深入研究，将复合命题中的命题常项换为命题变项，称这样的式子为合式公式，合式公式由命题变项、联接词、括号组成。
\ITEMS
	\item 单个命题变项是合式公式。
	\item 如果 $A$ 是合式公式，则 $\neg A$ 也是合式公式。
	\item 如果 $A,B$ 是合式公式，则用联接词联接 $A,B$ 也是合式公式。
	\item 只有有限次地应用 1,2,3 组成的式子才是合式公式。
\ITEME
在命题逻辑中，称合式公式为命题公式。

\subsubsection{赋值}

设 $A$ 为一个命题公式，$p_1,p_2,\cdots,p_n$ 为出现在 $A$ 中的所有命题变项。给 $p_1,p_2,\cdots,p_n$  指定一组真值，称为对 $A$ 的一个赋值或解释。若指定的一组值使 $A$ 的值为真，则称这组值为 $A$ 的成真赋值，反之为成假赋值。

\subsubsection{命题公式的分类}

设 $A$ 为一个命题公式
\ITEMS
	\item 若 $A$ 在所有赋值下取之均为真，则称 $A$ 为重言式或永真式。
	\item 若 $A$ 在所有赋值下取之均为假，则称 $A$ 为矛盾式或永假式。
	\item 若 $A$ 至少存在一组成真赋值，则称 $A$ 是可满足式。
\ITEME

\subsection{等值演算}

设 $A,B$ 是两个命题公式，若等价式 $A \leftrightarrow B$ 是重言式，则称 $A,B$ 是等值的，记作 $A\Leftrightarrow B$ 。

\begin{remark}
在讨论逻辑学时 $\leftrightarrow,\Leftrightarrow$ 两个符号的含义不同，$\Leftrightarrow$ 并不是联接词，但在分析学时我们会使用 $\Leftrightarrow$。
\end{remark}

根据已知的等值式，推演出与给定公式等值的公式的过程称为等值演算

\begin{theorem}
设 $\Phi(A)$ 是含命题公式 $A$ 的命题公式，$\Phi(B)$ 是用命题公式 $B$ 置换了 $\Phi(A)$ 中的 $A$ 之后得到的命题公式。如果 $A\Leftrightarrow B$ ，则 $\Phi(A)\Leftrightarrow\Phi(B)$。
\end{theorem}

\section{范式}

\subsection{范式}

仅由有限个命题变项或其否定构成的析取式称为简单析取式，仅由有限个命题变项或其否定构成的合取式称为简单合取式。

仅由有限个简单合取式构成的析取式称为析取范式，仅由有限个简单析取式构成的合取式称为合取范式。

\begin{theorem}[范式存在定理]
任意命题公式都存在与之等值的析取范式和合取范式。
\end{theorem}

\subsection{主范式}

设有 $n$ 个命题变项，若在简单合取式中每个命题变项与其否定有且仅有一个出现一次，则这样的简单合取式称为极小项。

如果公式 $A$ 的析取范式中的简单合取式全是极小项，则称该析取范式为 $A$ 的主析取范式。

\begin{theorem}
任意命题公式都有唯一的主析取范式。
\end{theorem}

设有 $n$ 个命题变项，若在简单析取式中每个命题变项与其否定有且仅有一个出现一次，则这样的简单析取式称为极大项。

如果公式 $A$ 的合取范式中的简单析取式全是极大项，则称该合取范式为 $A$ 的主合取范式。

\begin{theorem}
任意命题公式都有唯一的主合取范式。
\end{theorem}

\section{一阶逻辑基本概念}

\subsection{个体词}

个体词是指可以独立存在的个体，可以是具体的事物，也可以是抽象的概念。表示具体的或特定个体的词称为个体常项，表示抽象的或泛指的个体的词称为个体变项。个体变项的取值范围称为个体域。

\subsection{谓词}

谓词是用来刻画个体词的性质或个体词之间关系的词。表示具体性质和关系的谓词称为谓词常项，表示抽象的或泛指的谓词称为谓词变项。谓词中包含的个体词数称为元数，称包含 $n$ 个个体词的谓词为 $n$ 元谓词。

\subsection{量词}
量词指的是表示数量的词，有全称量词和存在量词两种。符号 $\forall$ 表示全称量词，读作“任意”，符号取自英文词汇 Arbitrary 的首字母。$\forall$ 用于表达某指定范围内“所有”、“任意”或“每一个”对象都具备某种性质。符号 $\exists$ 表示存在量词，读作“存在”，符号取自英文词汇 Exists 的首字母。$\exists$ 用于表达某指定范围内“至少存在一个”对象具备某种性质。

特别的，记 $\exists !$ 为唯一存在量词，用于表达某指定范围内“唯一存在一个”对象具备某种性质。


% ================================================================================== 
% 									集 合 论
% ================================================================================== 
\chapter{集合论}

集合论（set theory），是研究集合（由一堆抽象对象构成的整体）的数学理论，包含集合和元素（或称为成员）、关系等最基本数学概念。集合论、命题逻辑与谓词逻辑共同构成了数学的公理化基础，以未定义的“集合”与“集合成员”等术语来形式化地建构数学对象。

德国数学家格奥尔格·康托尔（Cantor）于十九世纪七十年代提出了朴素集合论，并从此展开了对集合论的研究。

\section{集合及其基本运算}

\subsection{集合及其基本计算}

\subsubsection{康托尔集合论}

康托尔集合论，也被称为“朴素”集合论，我们将其基本前提总结为以下几点

\ITEMS
	\item 互异性 \quad 集合可有任何有区别的对象组成。
	\item 无序性 \quad 集合由其组成对象整体唯一确定。
	\item 确定性 \quad 任何性质都确定一个具有该性质的对象的集合。
\ITEME

\subsubsection{包含关系}

我们将组成一个集合的对象称为该集合的元素，元素 $x$ 属于集合 $X$ 可以简记为
\[
	x\in X\ \mbox{或}\ X\ni x 
\]
其否命题为
\[
	x\notin X\ \mbox{或}\ X \not\ni x
\]
\subsubsection{简单的集合运算}

\ITEMS
	\item 集合$A$与$B$的并集是
	\[
		A\cup B = \{x\in U\mid (x\in A) \vee (x\in B)\}
	\]
	\item 集合$A$与$B$的交集是
	\[
		A\cap B = \{x\in U\mid (x\in A) \land (x\in B)\}
	\]
	% 在书中补集是作为差集的延伸提出的，但我们往往把补集视为基础运算的组合
	\item 集合$A$在全集$U$上的补集是
	\[
		A^c = C_U A= \{x\in U\mid (x\notin A) \}
	\]
	\item 集合$A$与$B$的差集是
	\[
		A\setminus B = \{x\in U\mid (x\in A) \land (x\notin B)\}
	\]
	应用以上基本运算，我们可以验证以下关系式（德摩根法则 A. de Morgan）
	\[
		(A\cup B)^c = A^c\cap B^c
	\]
	\[
		(A\cap B)^c = A^c\cup B^c
	\]
	\item 集合的直积（笛卡尔积）
	\[
		X\times Y=\{(x,y)\mid (x\in X)\land(y\in Y)\}
	\]
	其中 $(x,y)$ 为一个序偶，序偶由两个对象按照特定先后顺序排列构成的组合。
\ITEME

\section{映射与函数}

\subsection{映射}
设有集合 $X,Y$，如果 $X$ 的每一个元素 $x$ 都按照某种规律 $f$ 与 $Y$ 的唯一元素相对应，我们称其为一个映射，通常用以下符号表示映射
\[
	f:X\rightarrow Y
\]
对于 $\forall A\subseteq X$，我们称 $f(A)=\{f(x):x\in A\}$ 为 $A$ 的像；对于 $\forall B\subseteq Y$ ，我们称 $f^{-1}(B)=\{x\in X:f(x)\in B\} $ 为 $B$ 的原像。

特别的，我们称非空的数集之间的映射为函数。

\subsubsection{限制与延拓}

设有集合 $X,Y$，有映射 $f:X\to Y$，对于 $A \subseteq X$ 构建如下映射
\[
	f|_A : X\to Y
\]
映射 $f|_A$ 与 $f$ 在集合 $A$ 下相等，即在每个元素 $x\in A$ 上映射所得的像相等。我们称 $f|_A$ 为 $f$ 在 $A$ 上的限制，而相对于 $f|_A$ 称 $f$ 为 $f|_A$ 在 $X$ 上的延拓。

\subsubsection{映射的分类}
满足不同元素具有不同的像的映射称为单射，即
\[
	x_1 \ne x_2 \Rightarrow f(x_1) \ne f(x_2)
\]

满足映射的像等于集合 $Y$ 的映射称为满射，即
\[
	f(X)=Y
\]

即是单射又是满射的映射的映射是双射，即
\[
\begin{aligned}
	&f(X)=Y\\
	&(f(x_1) = f(x_2)) \Rightarrow(x_1=x_2)
\end{aligned}
\]
\subsubsection{逆映射}

当映射是双射的时候，存在$f:X\rightarrow Y$的逆映射，记为
\[
	f^{-1}:Y\rightarrow X
\]
\subsubsection{映射的复合运算}

如果在映射 $f:X\to Y$ 和 $g:Y\to Z$ 中，构建一个新的映射
\[
\begin{aligned}
	g\circ f:X&\to Z\\
	x &\mapsto g(f(x))
\end{aligned}
\]
这样构建出的映射 $g\circ f$ 称为映射 $f$ 与 $g$ 的复合映射。

\subsection{关系}

关系是两个集合笛卡尔积的子集。如果有两个集合 $A$ 和 $B$，它们之间的二元关系 $\mathcal{R}$ 是集合 $A \times B$ 的子集。若元素 $a \in A$ 与 $b \in B$ 存在关系，则记作 $a\mathcal{R}b$ 或 $(a, b) \in \mathcal{R}$。对于集合 $X$，如果 $\mathcal{R} \subseteq X^2$，就说在 $X$ 上给定了二元关系 $\mathcal{R}$。

\subsubsection{等价关系}

对于一非空集合 $X$，若在 $X$ 上给定了二元关系 $\mathcal{R}$ 满足以下条件
\ITEMS
	\item 自反性\quad $\forall x\in X \Rightarrow x\mathcal{R}x$
	\item 对称性\quad 若 $x,y\in X , x\mathcal{R}y$ 则 $y\mathcal{R}x$
	\item 传递性\quad 若 $x,y,z\in X , x\mathcal{R}y, y\mathcal{R}z$ 则 $x\mathcal{R}z$
\ITEME 
则称该二元关系为等价关系，记作 $x\sim y$。

% 偏序关系
% 函数关系

\section{公理化集合论}

\subsection{罗素悖论与第三次数学危机}

罗素悖论设集合$A=\{x\mid x\in A\}$，那么我们可以假设$A\in A\Rightarrow A\notin A$，同样$A\notin A \Rightarrow A\in A$，我们得出了截然不同的两种结果。

这意味着现代数学基础的集合论存在瑕疵，数学家们随即对集合论进行了严格的约束，从而规避了逻辑悖论，但第三次数学危机仍并未被“完美”解决，之后的哥德尔不完备定理表明任何系统都无法做到既完全又自洽，从此促使数学从追求绝对真理转变为现代公理化体系。

\subsection{公理化集合论}

实际上公理化集合论是一个没有数学美感的产物，就像东拼西凑出的支柱，主要为防止如罗素悖论这种矛盾的产生，其中我们最常用的是策梅洛-弗伦克尔公理系统（Zermelo-Fraenkel），一般简称为 ZF 集合论。
\ITEMS
	\item 外延公理
	
	集合$A$与集合$B$相等，当且仅当它们所具有的各元素是相同的，即
	\[
		A=B\iff A\subset B \land A\supset B
	\]
	\item 分离公理
	
	任何集合$A$和性质$P$都对应一个集合，其元素是且仅是$A$中具有性质$P$的各元素，即
	\[
		\{x\in A: P(x)\}\subset A
	\]
	
	\item 并集公理
	
	分离公理为集合运算的基础中的交集和补集提供了公理支持，但是却无法支持并集运算，那么我们提出对于集合的任何集合$M$，存在一个集合$M$的并集的集合$\cup M$，其元素是且仅是$M$的元素所包含的那些元素。
	\item 配对公理
	
	对于任何集合 $X$ 和 $Y$ ，存在一个集合 $Z$ ，其元素仅为 $X$ 和 $Y$。
	
	\item 空集公理
	
	存在一个不包含任何元素的集合 $\varnothing$，由分离定理从任意集合都可以分离出空子集，并根据外延定理可知所有空子集都相等。
	
	\item 幂集公理（子集之集）
	
	对于任何集合 $X$，存在一个集合 $\mathcal{P}$，其元素是且仅是 $X$ 的各子集。
	\item 无穷公理
	
	我们首先需要引入后继集的概念，定义 $X^+=X\cup \{X\}$ 为后继集，此外，如果一个集合包含空集以及自身任何一个元素的后继集，我们称该集合为归纳集。
	
	而无穷公理也就是归纳集存在。
	
	我们在无穷公理的基础上再通过之前所用的公理就可以建立自然数集 $\mathbb{N}_0$ 的模型（冯诺伊曼表示法），对自然数作如下表示。
	\[
		\begin{aligned}
			0:\quad & \varnothing\\
			1:\quad & \varnothing^+=\{\varnothing\}\\
			2:\quad & \{\varnothing\}^+=\{\{\varnothing\}\}\\
			\vdots &
		\end{aligned}
	\]
	\item 替换公理
	 
	设$\mathcal{F}(x,y)$是以下命题（公式）：对于集合$X$中的任何元素$x_0$存在唯一的对象$y_0$，使得$\mathcal{F}(x_0,y_0)$成立，那么满足该条件的对象$y$组成一个集合。
	
	\item * 选择公理（策梅洛公理）
	
	对于任意互不相交的非空集合组成的集合族，存在集合$C$，使得对于该集合族中的任何集合$X$，集合$X\cap C$只有一个元素组成。
	
	ZF 集合论不包含选择公理，因为选择公理存在一些争议，例如承认此公理会引出分球怪论，但由于选择公理可以为证明提供很多便利，我们有 ZFC 集合论，也就是在 ZF 集合论再加上选择公理。
\ITEME

\section{实数集公理系统与一般性质}

\subsection{实数集的定义}

如果以下四组条件成立，则集合 $\R$ 称为实数集，其元素称为实数，这些条件称为实数公理系统。

\ITEMS
	\item 加法公理
	
	定义如下映射（加法运算）
	\[
		+:\R\times \R\rightarrow\R
	\]
	
	1$_+$. 中性元素 $0$ 存在;，并且对任何的 $x\in \R$，
	\[
		x+0 = 0+x = x
	\]
	
	2$_+$. 对于任何元素 $x\in\R$，元素 $-x\in\R$ 一定存在，称为 $x$ 的相反元素，满足
	\[
		x+(-x)=(-x)+x=0
	\]
	
	3$_+$. 运算满足结合律，即对 $\R$ 的任何元素 $x,y,z$ 满足
	\[
		x+(y+z) = (x+y)+z
	\]
	
	4$_+$. 运算满足交换，即对 $\R$ 的任何元素 $x,y$ 满足
	\[
		x + y = y + x
	\]
	
	如果在某集合 $G$ 上定义了满足前三条公理的运算，我们就说在 $G$ 上给定了群结构，即 $G$ 是群，如果成该运算为加法，就称这个群为加法群，如果除此之外还知道该运算满足交换律，即条件 4$_+$ 成立，就称这个群为交换群或阿贝尔群。
	
	\item 乘法公理
	
	定义如下映射（乘法运算）
	\[
		\bullet:\R\times \R\rightarrow\R
	\]
	
	1$_\bullet$. 中性元素 $1\in\R\setminus0$ 存在;，并且对任何的 $x\in \R$，
	\[
		x\cdot 1 = 1\cdot x = x
	\]
	
	2$_\bullet$. 对于任何元素 $x\in\R\setminus0$，元素 $-x\in\R$ 一定存在，称为 $x$ 的逆元素，满足
	\[
		x+(-x)=(-x)+x=0
	\]
	
	3$_\bullet$. 运算满足结合律，即对 $\R$ 的任何元素 $x,y,z$ 满足
	\[
		x\cdot(y\cdot z) = (x\cdot y)\cdot z
	\]
	
	4$_\bullet$. 运算满足交换，即对 $\R$ 的任何元素 $x,y$ 满足
	\[
		x \cdot y = y \cdot x
	\]
	我们指出集合 $\R\setminus0$ 相对于乘法运算是（乘法）群。
	
	乘法相对于加法满足分配律，即 $\forall x,y,z \in \R$，
	\[
		(x+y)\cdot z = x \cdot z + y \cdot z
	\]
	
	如果某集合 $G$ 上定义了满足上述全部公理的两种运算，则 $G$ 称为代数域。
	
	\item 序公理
	
	对于$\R$的元素$x$，$y$，可以存在$x\le y$是否成立，这时以下条件成立
	
	0$_\le$. $\forall x\in\R (x\le x)$
	
	1$_\le$. $(x\le y)\land(y\le x)\Rightarrow(x=y)$
	
	2$_\le$. $(x\le y)\land(y\le z)\Rightarrow(x\le z)$
	
	3$_\le$. $\forall x\in\R\forall y\in\R(x\le y)\lor(y\le x)$
	
	如果一个集合的某些元素之间的关系满足公理 0$_\le$，1$_\le$，2$_\le$，该集合就称为偏序集，如果除此之外还满足 3$_\le$，即集合的任何两个元素都是可比较的，该集合就称为线性序集。
	
	
	\item 完备性公理
	
	如果 $X$ 和 $Y$ 是 $\R$ 的非空子集，并对于任何 $x\in X$，$y\in Y$ 有 $x \le y$，则存在 $c\in\R$，使得对于任何元素 $x\in X$，$y\in Y$ 有 $x\le c\le y$。
\ITEME

% 实数集代数性质

\subsection{完备性公理与数集的上确界（下确界）的存在性}

\subsubsection{上下界}

设集合$X\subset \R$，如果数$c\in\R$存在，使得对于任何$x\in X$都有$x\le c$，就说集合$X$是上有界集，这时称$c$为$X$的上界。

同理可以定义下界，既有上界又有下界的集合称为有界集。

\subsubsection{最值}

设 $a$ 是集合 $X\subset\R$ 的元素，如果对于任何元素 $x\in X$ 都有 $x\le a$ 则元素 $a$ 称为集合 $X$ 的最大元素（最大值）记作 $\max X$ 或 $\max\limits_{x\in X} x$，而最小值记作 $\min X$ 或 $\min\limits_{x\in X} x$。
\[
\begin{aligned}
	(a = \max X) := (a\in X \land \forall x \in X (x\le a))\\
	(a = \min X) := (a\in X \land \forall x \in X (a\le x))	
\end{aligned}
\]

\subsubsection{上下确界}

集合 $X\subset\R$ 的上界中的最小者称为集合 $X$ 的上确界，同理定义下确界如下
\[
\begin{aligned}
	(s = \sup X) &:= (\forall x\in X\,(x\le s))\land (\forall s' < s\, \exists x' \in X (s'< x'))\\
	(i = \inf X) &:= (\forall x\in X\;(i\le x))\land (\forall i' > i\; \exists x' \in X (x'< i'))
\end{aligned}
\]
也有等价定义
\[
\begin{aligned}
	\sup X  := \min\{c\in\R\mid\forall x\in X(x\le c)\}\\
	\inf X  := \max\{c\in\R\mid\forall x\in X(c\le x)\}
\end{aligned}
\]

\begin{lemma}[上确界原理]
实数集的任何有上界非空子集有唯一的上确界。
\end{lemma}
\begin{proof}
	设 $X \subseteq \R$ 是给定的子集，而 $Y = \{y\in\R | \forall x \in X (x \le y)\}$ 是 $X$ 的上界的集合，依条件 $X\neq \varnothing, Y\neq \varnothing$。于是，根据完备性公理，存在 $c \in \R$，使得 $\forall x \in X,\forall y \in Y (x\le c\le y)$，因此数 $c$ 既是 $X$ 的上界，也是 $Y$ 的下届，所以有 $c = \min Y = \max X$。
\end{proof}

\section{实数类与计算}

\subsection{自然数集}

\subsubsection{皮亚诺算数公理}
如果以下条件成立，则集合 $\N$ 称为自然数集，其元素称为自然数。
\ITEMS
	\item $0$ 是自然数。
	\item 对于集合 $N$ 中的任意元素 $n$，其后继 $n'$ 也是 $N$ 中的一个明确且唯一的自然数。
	\item $0$ 不是任何自然数的后继数。
	\item 对于每个自然数 $b,c$， $b=c$ 当且仅当 $b' = c'$。
	\item 任意关于自然数的命题，如果证明：它对自然数 $0$ 是真的，且假定它对自然数 $a$ 为真时，可以证明对 $a'$ 也真，那么，命题对所有自然数都真。
\ITEME
\begin{remark}
	其中条件 5 保证了数学归纳原理的正确性，是数学证明的基本工具。
\end{remark}

\subsection{有理数与无理数}

\subsubsection{整数集}


\subsubsection{有理数集}

\subsubsection{无理数集}

\subsection{数轴}

\subsubsection{区间}

\subsubsection{邻域}

包含点 $x \in \R$ 的开区间称为该点的邻域，特别的，当 $\delta > 0$ 时，开区间 $(x - \delta,x+\delta)$ 称为点 $x$ 的 $\delta$ 邻域。

\subsubsection{模}



\section{关于实数集完备性的一些基本引理}
% 1
\subsection{闭区间套引理（柯西-康托尔原理）}
\subsubsection{序列}

我们将以自然数集为定义域的映射 $f:\mathbb{N}\rightarrow X$ 称为序列。

\subsubsection{集合套}

设$X_1,X_2,\cdots,X_n,\cdots$是某些集合的序列，如果$X_1\supset X_2\supset \cdots \supset X_n \supset\cdots$，我们称之为集合套序列。

\begin{lemma}[闭区间套引理]
	对于任何闭区间套序列 $I_1\supset I_2\supset \cdots \supset I_n \supset\cdots$，可以找到一点 $c\in\R$ 属于所有这些闭区间。
	
	如果除此之外还已知，对于 $\forall \varepsilon > 0$，在序列中可以找到长度 $|I_k| < \varepsilon$ 的闭区间，则 $c$ 是所有闭区间的唯一公共点。
\end{lemma}
% 使用完备性公理证明
\begin{proof}
首先	指出，对于上述序列中的任何两个闭区间 $I_m = [a_m,b_m],I_n = [a_n,b_n]$ 必有 $a_m \le b_n$。因此，对于数集 $A =\{a_m | m\in \N\}, B= \{b_n| n\in \N\}$，根据完备性公理，可以找到数 $c\in\R$ 使得对于 $\forall a_m\in A,\forall b_n\in B$ 有 $a_m\le c\le b_n$。特别的，对于 $\forall n\in N$ 都有 $a_n\le c \le b_n$，即点 $c$ 属于所有闭区间 $I_n$。

现在，假设有两个不同的有相似性质的点 $c_1,c_2$，假如 $c_1 < c_2$，则对于任何 $n\in\N$ 均有， $a_n \le c_1 < c_2 \le b_n$ 从而 $0 < c_1 - c_2 < b_n - a_n$。于是，上述序列中的每一个闭区间的长度不可能小于 $c_2 - c_1 > 0$，这意味着，如果在该序列中有长度任意小的闭区间则公共点是唯一的。 
\end{proof}


\subsection{有限覆盖引理（博雷尔-勒贝格原理）}

\subsubsection{覆盖}

设 $S=\{X\}$ 是由一组集合 $X$ 构成的集合族，如果 $Y\subset \bigcup\limits_{X\in S}X$，那么我们说 $S$ 覆盖集合 $Y$。

集合族 $S = {X}$ 的子集也是一个集合组，称为 $S$ 的子族。因此，集合族的子族本身是同一类型的集合族。

\begin{lemma}[博雷尔-勒贝格原理]

在覆盖一个闭区间的任何开区间族中都有覆盖该闭区间的有限子族。
\end{lemma}
\begin{proof}

\end{proof}


\subsection{极限点引理（波尔查诺-魏尔斯塔拉斯原理）}

\subsubsection{极限点}
如果点 $p\in\R$ 的任何邻域都包含集合 $X\subset\R$ 的一个无穷子集，点 $p$ 就称为集合 $X$ 的极限点。

\begin{lemma}[波尔查诺-魏尔斯塔拉斯原理]
任何无穷有界数集至少有一个极限点。
\end{lemma}
\begin{proof}

\end{proof}

\section{集合的势}

\subsection{集合的势}

粗略地讲，任一集合 $A$ 的基数或者势，记为 $|A|$，是指 $A$ 所含元素的个数。我们称自然数集
\[
	\mathbb{N}=\{1,2,\cdots\}
\]
的基数为可数无穷，记为 $\aleph_0$。

在比较任意两个集合 $A$ 和 $B$ 的基数时，我们有下面的定义：

\ITEMS
	\item 我们称 $\lvert A \rvert \le \lvert B \rvert$，如果存在从 $A$ 到 $B$ 的单射。
	\item 我们称 $\lvert A \rvert = \lvert B \rvert$，如果$\lvert A \rvert \le \lvert B \rvert$ 且 $\lvert B \rvert \le \lvert A \rvert$，也就是说，存在从 $A$ 到 $B$ 的双射。
\ITEME


\subsubsection{幂集}
称任一集合$A$的所有子集构成的集合为集合$A$的幂集，记为$P(A)$。不难看出
\[
\lvert P(A) \rvert = 2^{\lvert A \rvert}
\]

\begin{theorem}[无最大基数定理]
对于任一非空集合 $A$，我们总有 $\lvert A \rvert<\lvert P(A) \rvert$。
\end{theorem}

\begin{proof}

我们先证明 $\lvert A \rvert \le \lvert P(A) \rvert$：\\
定义映射 $g:A\rightarrow P(A)$ 如下：
\[ 
	g(x)=\{x\}, \forall x \in A
\] 
不难看出，$g$ 是单射。

下面我们用反证法证明 $\lvert A \rvert \ne \lvert P(A) \rvert$:
若相等，则存在双射
\[
	f:A\rightarrow P(A)
\]
定义一个$A$的子集合$B$如下：
\[
B = \{x\in A: x\notin f(x)\}
\]
由于 $f$ 是双射，所以必存在（唯一）元素 $y\in A$ 使得
\[
f(y) = B
\]
考虑 $y$ 与 $B$ 的关系，发现无论如何都会产生矛盾：
\ITEMS
	\item 如果 $y\in B=f(y)$, 由 $B$ 的定义，我们会得到 $y\notin B$；
	\item 如果 $y\notin B=f(y)$, 由 $B$ 的定义，我们会得到 $y\in B$。
\ITEME

我们记 $\aleph_n=2^{\aleph_{n-1}}, n=1, 2, \cdots$。由上述定理，显然我们有 
\[
	\aleph_0 < \aleph_1 < \aleph_2 <\cdots< 
\]
\end{proof}

\subsubsection{关于基数的几个基本定理}
\ITEMS
	\item 对任意集合 $A$ 和 $B$，若 $A\subseteq B$，则 $\lvert A \rvert \le \lvert B \rvert$。
	\item 设 $A$ 是可数集（有限集或可数无穷集），$B$ 是无穷集，则 $\lvert A\cup B \rvert= \lvert B \rvert$。
	\item 设集合列 $\{A_n\}_{n=1}^{\infty}$ 中每一 $A_n$ 的基数都是 $\aleph_0$，则 $\lvert \bigcup_{n=1}^\infty A_n \rvert=\aleph_0 $。也就是说，可数个可数集的并集还是可数集。
\ITEME


\subsubsection{实数集的基数}
实数集 $\R$ 的基数一般记为 $c$。

\begin{proof} $c = \aleph_1$

我们先证明 $\lvert \R \rvert=\lvert (0,1]\rvert $，分三步进行：\\

1. $\lvert \R\rvert =\lvert (-1,1)\rvert $: 不难看出映射  $f:(-1,1)\rightarrow \R$ 满足 $f(x)=\frac {x}{1-x^2}$ 是双射。\\

2. $\lvert (-1,1)\rvert=\lvert (0,1)\rvert $: 不难看出映射  $f:(-1,1)\rightarrow (0,1) $ 满足 $f(x)=\frac {x+1}{2}$ 是双射。\\

3. $\lvert (0,1)\rvert=\lvert (0,1]\rvert $： 因为 $(0,1)\subseteq (0,1]\subseteq (-1,1)$。\\

下面我们证明 $\lvert (0,1]\rvert=2^{\aleph_0}$: \\
1. 把 $(0,1]$ 中的任一实数 $x$ 用二进制表示：更具体地讲，就是找到每项取值是 $0$ 或 $1$ 的无穷数列 $\{a_n\}_{n \ge 1}$ 满足在二进制下 $x=0.a_1 a_2\cdots a_n \cdots$，其中无穷多个 $a_n$ 是 $1$ 。\\
\ITEMS 
	\item 第一步: 把区间 $(0,1]$ 二等分为两半开闭子区间：若 $x$ 属于靠左的（称为第一个，后同）子区间 $(0,\frac{1}{2}]$，则 $a_1=0$；若 $x$ 属于靠右的（称为第二个，后同）子区间，则 $a_1=1$；
	
	\item 第二步：把 $x$ 所属的第一步中的小区间二等分为两半开闭子区间:  若 $x$ 属于第一个子区间，则 $a_2=0$；若 $x$ 属于第二个子区间，则 $a_2=1$；
	\item $\vdots$ 
	\item 第 $n$ 步： 把 $x$ 所属的上一步中的小区间二等分为两半开闭子区间:  若 $x$ 属于第一个子区间，则 $a_n=0$；若 $x$ 属于第二个子区间，则 $a_n=1$；
	\item $\vdots$
\ITEME
    注意到上述 $a_n$ 中，有无穷多个是 $1$ （为什么？）。这样我们就建立了 $(0,1]$ 和集合 $$S=\{(a_1, a_2,\cdots, a_n, \cdots): \mbox{每个 $a_n$ 是 $0$ 或 $1$ 且 其中无穷多个 $a_n$ 是 $1$} \}$$ 的一一对应关系。\\
2 决定 $\lvert S\rvert$：令 $$T=\{(a_1, a_2,\cdots, a_n, \cdots): \mbox{每个 $a_n$ 是 $0$ 或 $1$ 且 其中只有限多个 $a_n$ 是 $1$} \}$$ 和 $$ R=\{(a_1, a_2,\cdots, a_n, \cdots): \mbox{每个 $a_n$ 是 $0$ 或 $1$}\}$$ 
不难看出 $\lvert R\rvert=2^{\aleph_0}, \lvert T\rvert=\aleph_0$ （为什么？）且 $S\cup T= R$。由上面关于基数的基本定理 $2$，我们有 $$\lvert S\rvert=\lvert S\cup T\rvert =\lvert R\rvert=2^{\aleph_0}=\aleph_1$$
\end{proof}

\subsubsection{连续统假设}

不存在一个无穷集合，其基数严格介于自然数集的基数 $\aleph_0$ 与实数集的基数 $\aleph_1$（即连续统的基数）之间。

不存在一个无穷集合，其基数严格介于 $\aleph_n$ 与 $\aleph_{n+1}, n=0, 1,2, \cdots $ 之间。

\section{数学史：康托尔}

% ================================================================================== 
% 									极 限 论
% ================================================================================== 
\chapter{极限论}

\section{数列极限}

% 介绍数列的概念
\subsection{数列极限的定义}

% 特殊的函数极限
设$\{x_n\}$为一数列，如果存在一常数$a$，对于任意给定的正数$\varepsilon$，总存在正整数$N$，使得任意$n > N$
\[
	|x_n-a|<\varepsilon
\]
都成立，那么就称常数$a$是数列$\{x_n\}$的极限，或者称数列$\{x_n\}$收敛于$a$，记为
\[
	\lim_{n \to \infty}x_n = a 
\]
如果不存在这样的常数$a$，就说数列$\{x_n\}$没有极限，或者数列$\{x_n\}$是发散的

\[
\forall \varepsilon > 0, \exists N \in \mathbb{N}_+\ s.t.\ n>N \Rightarrow |x_n - A|<\varepsilon
\]

\subsection{收敛数列的性质}

% 唯一性 保号性 保序性 迫敛性 任意子列收敛

\ITEMS
\item 唯一性

如果数列收敛，那么它的极限唯一

\item 保号性

如果数列极限为$a$，且$a>0$，或$a<0$，那么存在正整数$N$，对于$\forall n> N$，都有$x_n>0$

\item 有界性

如果数列收敛，那么数列一定有界

\item 子列收敛性

如果数列极限为$a$，那么它的任意子数列也收敛也为$a$
\ITEME

\subsection{数列极限的存在问题}

\subsubsection{柯西准则}
% 因为定义中的A往往是未知的，所用极限不便于应用
% 三次数学危机，第二次数学危机，无穷小量的引入
对于数列$\{x_n\}$收敛的充要条件是，对于任意给定的正数$\varepsilon$，存在正整数$N$，使得当$n,m > N$时，满足
\[
	|x_n-x_m| < \varepsilon
\]
% 体现数列逐渐稠密
我们称满足以上条件的数列为基本数列（柯西数列）
\[
\forall \varepsilon \in \mathbb{N}_+ ,\ \exists N,\ n,m>N \implies |x_n-x_m| < \varepsilon
\]

\subsubsection{单调数列极限存在准则（魏尔斯特拉斯定理）}

单调不减数列的极限存在的充要条件是有上界

\subsubsection{自然常数}

\subsubsection{子列与数列的部分极限}

\subsection{级数初步}

\subsubsection{级数的柯西准则}



\section{函数极限}

\subsection{函数极限的定义}

设函数$f(x)$在点$x_0$的某一{\bf 空心邻域}内有定义，如果存在常数$A$，对于任意给定的正数$\varepsilon$，总存在正数$\delta$，使得当$x$满足不等式$0 < |x-x_0| <  \delta$时对应的函数值都满足
\[
	|f(x)-A|<\varepsilon
\]
那么就称$A$为函数$f(x)$当$x\to x_0$时的极限，记作
\[
\lim_{x\to x_0} f(x)= A
\]


符号化
\[
\forall \varepsilon > 0 , \exists \delta > 0\ s.t.\  0 < |x-x_0| <  \delta \implies |f(x)-A| < \varepsilon\
\]

% 七种极限定义
% 左右单侧极限
% 无穷极限
% 极限趋近于无穷（极限不存在）


\subsection{函数极限的性质}
% 这一部分还不完善	<-------------------------- HERE

\subsection{极限的一般定义}

\subsection{函数极限的存在问题}

\section{函数的连续性}

\subsection{连续性的定义}

\subsection{连续函数的性质}


% 在连续性的定义上不必是空心邻域
$f(x)$在点$x_0$点连续 
\[
\lim_{x\to x_0} f(x) = f(x_0) 
\]

% 区间上连续

% 一致连续性（扩展）

%\[
%\lim_{x \to x_0}f(x) = A, \mbox{求}\lim_{x \to x_0} e^{f(x)}
%\]

% 幂指函数极限

%\[
%f(x)^{g(x)} = e^{\ln f(x)^{g(x)}} =e^{g(x)\ln f(x)}
%\]


\subsection{间断点}


\subsection{零值定理}
$f(x)$在$[a,b]$上连续，$f(a)f(b)<0$，则$\exists \xi \in (a,b)$使$\xi = 0$

% ================================================================================== 
% 									微 分 学
% ================================================================================== 
\chapter{微分学}
\section{导数与微分}
\subsection{可微函数}
% 这一部分还不完善	<-------------------------- HERE


\section{基本微分法则}

\section{微分学的基本定理}

\subsection{费马引理}

\subsubsection{局部极值点}



\subsection{罗尔定理}

假如函数\(f(x)\)在\([a,b]\)区间上连续，在\((a,b)\)区间上可微，且\(f(a)=f(b)\)，那么$\exists\, \xi\in(a,b)$使得$f'(\xi)=0$

\subsection{拉格朗日有限增量定理}

拉格朗日定理是罗尔定理的一般化情况，假如函数\(f(x)\)在\([a,b]\)区间上连续，在\((a,b)\)区间上可微，那么$\exists\, \xi\in(a,b)$使得
\[
	f'(\xi)=\frac{f(b)-f(a)}{b-a}
\]

\subsubsection{常函数检验法}
连续函数在闭区间$[a,b]$上为常数的充要条件是其导数在闭区间$[a,b]$上恒为零

\subsection{柯西有限增量定理}
设$x=f(t)$和$y=g(t)$是在闭区间$[a,b]$上连续且在开区间$(a,b)$上可微的函数，则存在点$\tau\in(a,b)$使得
\[
	f'(\tau)(g(b)-g(a)) = g'(\tau)(x(b)-x(a))
\]

\subsection{泰勒公式}

\subsubsection{泰勒多项式}
在数学的历史上我们经常会有用简单函数描述复杂函数的情况，所以我们希望能够用简单的多项式函数来描述所有的函数，假如有如下函数可以写成多项式之和的形式
\[
	P_n(x) = c_0 + c_1 x + c_2 x^2 + \cdots + c_n x^n
\]
对于该多项式我们可以用函数来表示多项式的系数
\[
	\begin{aligned}
		& P_n(0) = 0! c_0\\
		P_n'(x) = 1! c_1 + 2 c_2 x + \cdots + n c_n x^{n-1}\qquad & P_n'(0) = 1! c_1\\
		P_n''(x) = 2! c_2  + \cdots + n(n-1) c_n x^{n-2}\qquad & P_n''(0) = 2! c_2\\
		\vdots \qquad\qquad\qquad & \qquad\vdots\\
		P_n^{(n)}(x) = n! c_n\qquad & P_n^{(n)}(0) = n! c_n
	\end{aligned}
\]
函数可以重新写为
\[
	P_n(x) = \frac{P_n(0)}{0!} + \frac{P_n'(0)}{1!} x + \frac{P_n''(0)}{2!} x^2 + \cdots + \frac{P_n^{(n)}(0)}{n!} x^n 
\]
因为此多项式仅能用于描述$x=0$的邻域，我们还希望在该多项式的基础上进行拓展用于描述更多情况，我们用同样的方法可以写出以下多项式
\[
	\begin{aligned}
		P_n(x)=P_n(x_0;x) &= c_0 + c_1 (x - x_0) + c_2 (x - x_0)^2 + \cdots + c_n (x - x_0)^n\\
		&=\frac{P_n(0)}{0!} + \frac{P_n'(0)}{1!} (x - x_0) + \frac{P_n''(0)}{2!} (x - x_0)^2 + \cdots + \frac{P_n^{(n)}(0)}{n!} (x - x_0)^n 
	\end{aligned}
\]
如果给定了在$x_0$点具有所有前$n$阶导数的函数$f(x)$我们可以写出以下多项式
\[
	P_n(x_0;x) =\frac{f(0)}{0!} + \frac{f'(0)}{1!} (x - x_0) + \frac{f''(0)}{2!} (x - x_0)^2 + \cdots + \frac{f^{(n)}(0)}{n!} (x - x_0)^n 
\]
称之为函数$f(x)$的$n$阶泰勒多项式
\subsubsection{余项}



\section{微分与函数}

\section{复数与初等函数}

\section{微分的应用}

\section{原函数与不定积分}

% ================================================================================== 
% 									积 分 学
% ================================================================================== 
\chapter{积分学}

\section{积分的定义}

我们对于积分学的研究可以从两个角度开始，首先考虑如下物理问题。现有一物体做变加速度直线运动，试求其从 $a$ 时刻到 $b$ 时刻所经过路程（位移），其速度与时间图像如下
\begin{figure}[h]
\centering
\begin{tikzpicture}[scale=1.15]

\def\a{1.5}   
\def\b{4.5}   
\pgfmathsetmacro{\va}{0.025*\a^3 + 0.2*\a}
\pgfmathsetmacro{\vb}{0.025*\b^3 + 0.2*\b}
\draw[->, thick] (0,0) -- (5.8,0) node[below right] {$t$};
\draw[->, thick] (0,0) -- (0,5.8) node[above left] {$v$};

% 刻度
\foreach \x in {1,2,3,4,5}
    \draw (\x,0.15) -- (\x,-0.15) node[below] {\x};
\foreach \y in {1,2,3,4,5}
    \draw (-0.15,\y) -- (0.15,\y) node[left=8pt] {\y};
    
\draw[thick, black, smooth, samples=120, domain=0:5.5] 
    plot (\x,{0.025*\x^3 + 0.2*\x});

\fill[opacity=0.3] 
    plot[domain=\a:\b, samples=100] (\x,{0.025*\x^3 + 0.2*\x}) 
    -- (\b,0) -- (\a,0) -- cycle;
% 垂线
\draw[dashed, thick] (\a,0) -- (\a,\va);
\draw[dashed, thick] (\b,0) -- (\b,\vb);
\fill (\a,\va) circle (2pt);
\fill (\b,\vb) circle (2pt);
% 标注 
\node[below] at (\a, -0.2) {a};
\node[below] at (\b, -0.2) {b};
\node[above, font=\small] at (\a, \va + 0.4) {$v(a)$};
\node[above, font=\small] at (\b, \vb + 0.4) {$v(b)$};

\end{tikzpicture}
\end{figure}

显然该图像所围成的面积就是位移，我们可以将时间分割成 $n$ 个区间，其中第 $i$ 个区间大小为 $\Delta x_i$，记所用区间的最大值为 $\lambda = \max\limits_{1\le i\le n}\{\Delta x_i\}$ 。现在我们在每个区间上用矩形逼近函数图像，用 $\xi_i$ 表示在第 $i$ 个区间上的中间点，那么第 $i$ 个区间上的矩形面积应为 
\[
	v(\xi_i)\Delta t_i
\]
我们可以将总体的面积表示为
\[
	\sum_{i=1}^n v(\xi_i)\Delta t_i
\]
但这个面积仍不等于速度与时间图像的面积，为了使两者相等，要把时间分割成无穷个无穷小的区间，那么考虑对面积取如下极限
\[
	\lim_{\lambda\to0}\sum_{i=1}^n v(\xi_i)\Delta t_i
\]
同样的，在数学上，对于曲线函数 $f(x)$ 在区间 $[a,b]$ 上围成的曲边梯形，我们可以用相似的办法求得如下面积
\[
	\lim_{\lambda\to0}\sum_{i=1}^n f(\xi_i)\Delta x_i
\]
我们使用莱布尼兹的符号把这个极限记为
\[
	\int_a^b f(x)dx = \lim_{\lambda\to0}\sum_{i=1}^n f(\xi_i)\Delta x_i
\]
所以积分本质上是和式极限

\section{积分的性质}

\section{微积分}

\section{积分的应用}

\section{反常积分}


% ================================================================================== 
% 									拓 扑 学
% ================================================================================== 

\chapter{拓扑学}

% ================================================================================== 
% 								   线 性 代 数
% ================================================================================== 

\chapter{线性代数}

% ================================================================================== 
% 							  多 元 函 数 微 分 学
% ================================================================================== 

\chapter{多元函数微分学}

% ================================================================================== 
% 								  	重 积 分
% ================================================================================== 

\chapter{重积分}

% ================================================================================== 
% 							 曲 线 与 曲 面 积 分
% ================================================================================== 

\chapter{曲线与曲面积分}

% ================================================================================== 
% 									级 数 
% ================================================================================== 

\chapter{级数}

% ================================================================================== 
% 								 傅 立 叶 变 换
% ================================================================================== 

\chapter{傅立叶变换}

% ================================================================================== 
% 								  抽 象 代 数 
% ==================================================================================

\chapter{抽象代数}

% ================================================================================== 
% 							测 度 与 勒 贝 格 测 度 
% ================================================================================== 

\chapter{测度与勒贝格测度}

% ================================================================================== 
% 								勒 贝 格 积 分
% ================================================================================== 

\chapter{勒贝格积分}

% ================================================================================== 
% 								  概 率 论
% ================================================================================== 

\chapter{概率论}

% ================================================================================== 
% 								  随 机 过 程
% ================================================================================== 

\chapter{随机过程}

% 文档结束
\end{document}



